\chapter{Введение}

\section{Цель}  

Произвести настройку протокола DHCP на маршрутизаторе
R - DHCP. Настроить маршрутизатор R1 как ретранслятор. Рассмотреть процесс
получения IP – адреса клиентскими устройствам
\section{Задачи}

\begin{enumerate}
\item Интерфейс Gi1/0/1 маршрутизатора R1 соединить с интерфейсом Gi1/0/1
маршрутизатора R2.
\item Компьютер VPC1 подключить к интерфейсу Gi1/0/2 маршрутизатора R1.
\item Компьютер VPC2 подключить к интерфейсу Gi1/0/3 маршрутизатора R1.
\item Настроить имя маршрутизаторам и назначить IP - адреса на интерфейсы
согласно таблице 6.2.1.
\item Настроить DHCP сервер на маршрутизаторе R - DHCP:
\item Создать пул с именем DHCP-10
\item Выдавать из пула диапазон адресов 10.10.10.100 - 10.10.10.200
\item Выдавать адрес шлюза по умолчанию 10.10.10.1
\item Настроить время аренды адреса на 30 минут
\item Создать пул с именем DHCP-172
\item Выдавать из пула диапазон адресов 172.16.10.100 - 172.16.10.200
\item Выдавать адрес шлюза по умолчанию 172.16.10.1
\item Настроить helper-address 192.168.100.2 на интерфейсах gi1/0/2 и gi1/0/3
маршрутизатора R1.
\item Настроить маршрутизацию.
\item Получить IP - адреса на компьютеры VPC1 – VPC2 по протоколу DHCP.
\end{enumerate}
 
\begin{figure}[H]  
    \centering
    \includegraphics[width=0.75\textwidth]{top.png}
    \caption{Пример собранной сети}
\end{figure}

\begin{figure}[H] 
    \centering
    \includegraphics[width=0.75\textwidth]{table.png}
    \caption{Таблица с настройками}
\end{figure}

\endinput